\section{Time Series Data Analytics}

Time series analytics refers to the process of finding meaningful patterns and
insights from data that is collected over time. Unlike traditional data and
cross-sectional data, time series data exhibits temporal dependencies, meaning
that the value at time $t$ can be influenced by values at previous time
points $t-1, t-2, \ldots$.

\subsection{Why Analytics for Time Series Data?}

Raw data shows what is happening, while analysis reveals the context and the
underlying reasons. It allows us to identify trends, seasonal patterns, and
anomalies in the data, which can be crucial for making informed decisions in
various domains such as finance, IoT, weather forecasting, and healthcare.\cite{AI&ML:23}

\subsubsection{Analytics in our project}

AirScout stores a large amount of time-series data inserted by sensor boxes.
Unfortunately, this system is not fully secure, as somebody could reverse
engineer the boxes and insert false data if they have an AirScout account.
Therefore, we can use analytics to detect anomalies in the data, which could
indicate that somebody is inserting false data.

Several statistical techniques can be used for this, such as:

\begin{itemize}
    \item Mean and Standard Deviation Analysis: Calculate the mean and standard deviation
          of the time series data. Data points that fall outside a certain number of
          standard deviations from the mean can be flagged as anomalies.
    \item Z-Score: Builds on the concept of the mean and standard deviation, calculating
          the Z-score for each data point to identify how many standard deviations a
          point is from the mean. The formula for the Z-score is:
          \[ Z = \frac{(X - \mu)}{\sigma} \]
          where \(X\) is the data point, \(\mu\) is the mean, and \(\sigma\) is the
          standard deviation. Data points with a Z-score above a certain threshold (e.g.,
          $|Z| > 3$) can be considered anomalies.
    \item Moving Average: The moving average smooths out short-term fluctuations and
          highlights longer-term trends or cycles. Anomalies can be detected by comparing
          the actual data points to the moving average and flagging those that deviate
          significantly.
    \item Exponential Smoothing: This technique gives more weight to recent observations
          while still considering past data. Anomalies can be identified by comparing the
          actual values to the exponentially smoothed values. The formula for exponential
          smoothing is:
          \[ S_t = \alpha * X_t + (1 - \alpha) * S_{t-1} \]
          where \(S_t\) is the smoothed value at time \(t\), \(X_t\) is the actual value
          at time \(t\), and \(\alpha\) is the smoothing factor (0 < \(\alpha\) < 1). A
          higher \(\alpha\) gives more weight to recent observations. It's exponential
          because the weights decrease exponentially for older data points.
\end{itemize}~\cite{AI&ML:24}

Those techniques can be implemented in the database layer directly using SQL
queries.

\begin{lstlisting}[language=PGTSSQL, caption=Anomaly Detection using Z-Score in TimescaleDB]
SELECT *
FROM (
    SELECT *,
        AVG(val) OVER() as average,
        STDDEV(val) OVER() as std
    FROM measurements
) stats
WHERE ABS(val - average) > (3 * std);
\end{lstlisting}

This query calculates the average and standard deviation of the \texttt{val} column in
the \texttt{measurements} table and selects rows where the absolute difference between
the value and the average exceeds three times the standard deviation,
indicating potential anomalies.

\subsection{Performance issues with Analytics}

Analytics can be computationally intensive if the dataset is large and the
queries are complex. Regular databases like plain MySQL or key-value stores have to
scan through all rows or values, even to calculate simple statistics like the mean
or standard deviation. Time-series databases, on the other hand, are optimized
for such operations and can often perform them much faster due to their data
structure and indexing strategies. This thesis already mentioned compression
techniques of time-series databases (see Section~\ref{sec:compression}), which
can also help to speed up analytics queries, as less data has to be read from
disk. Additionally, time series databases often have built-in functions for
common time-series analytics tasks that are far more efficient because they use
precomputed aggregations and indexes on time-based data, and also employ
techniques like downsampling to reduce the amount of data that needs to be
processed for analytics queries. This can lead to significant performance
improvements, especially when dealing with large datasets and complex queries.
It should also be noted that built-in functions are often faster because they are
implemented in a lower-level language and optimized for performance, while
user-defined functions or external analytics tools may not be as efficient.

Taking the example above, the Z-score calculation requires a full scan of the
\texttt{meas} table to compute the window functions \texttt{AVG} and
\texttt{STDDEV} across the entire dataset. In a traditional row-oriented
relational database, this operation has a time complexity of $O(N)$, where $N$
is the number of rows. This often leads to significant I/O bottlenecks as the
system must load complete row pages even if only a single column is needed (see
Section~\ref{sec:benchmarking}).

In contrast, a time-series database leveraging columnar storage would only need
to read the specific blocks containing the \texttt{val} column, drastically
reducing disk I/O. Furthermore, by utilizing \textit{continuous aggregates} (or
materialized views), sufficient statistics such as counts, sums, and sums of
squares can be precomputed and incrementally updated as new data is
ingested~\cite{TimescaleDB:03}. This allows the query engine to derive the mean
and standard deviation from pre-aggregated bucket statistics rather than
scanning the raw hypertable. In practice, this changes the cost of the
statistical baseline from scaling with the number of raw measurements to
scaling with the number of aggregate buckets that must be read.

\begin{lstlisting}[language=PGTSSQL, caption=Creating a Continuous Aggregate in TimescaleDB]
CREATE MATERIALIZED VIEW measurements_stats
WITH (timescaledb.continuous) AS
SELECT time_bucket('1 hour', time) AS bucket,
    COUNT(*) as n,
    SUM(val) as sum_val,
    SUM(val * val) as sumsq_val
FROM measurements
GROUP BY bucket;
\end{lstlisting}

\begin{figure}[htbp]
    \centering
    \includegraphics[width=0.8\textwidth]{images/tigerdata-continuous-aggregate.png}
    \caption{Visualization of precomputed continuous aggregates over time}\label{fig:continuous_aggregates}
\end{figure}

Statistics are often accessed frequently, and each fetch should be fast and
should not be out of date. To avoid having to rescan every row every time, we
can precompute aggregates and store them in a separate table. For exact global
statistics over a time range, these precomputed bucket statistics still need to
be combined, but this is far cheaper than reading all raw measurements again.

\subsection{Batch Processing vs. Real-Time Processing}

When dealing with time-series data, analytics can be performed in two primary
paradigms: batch processing and real-time processing.

\textbf{Batch processing} involves collecting data over a period and processing it all at once in large blocks or `batches'. This means that analytical queries are executed on historical data, usually on a scheduled basis (e.g., nightly or weekly). Batch processing is highly efficient for complex, resource-intensive operations like training machine learning models or generating comprehensive monthly reports. However, it introduces latency, since the insights derived from the data are not immediately available.

\textbf{Real-time processing} (or stream processing), on the other hand, deals with data continuously as it arrives. Analytics are performed incrementally on the incoming data stream, allowing for near-instantaneous insights and actions. In the context of our project, AirScout, real-time processing is crucial for immediate anomaly detection. If a sensor box starts sending highly abnormal values, real-time analytics can trigger an alert immediately rather than waiting for the next batch cycle.

Choosing between batch and real-time processing depends on the use case. While
real-time processing provides lower latency, it is often more complex to
implement and requires specialized streaming engines.~\cite{BatchVsRealTime:01}
Many modern time-series setups use a combination of both architectures, where
real-time processing is used for critical alerting and fast views, while batch
processing serves for deep, long-term analytics.
\subsection{Machine Learning and Transformers for Time Series Data}

While statistical methods like Z-scores or moving averages are effective for
basic anomaly detection and linear trend analysis, they often fall short when
dealing with highly complex, non-linear, and multivariate time-series data.
This is where Machine Learning (ML), and more recently, Deep Learning models
based on the Transformer architecture come into play.\\\cite{AI&ML:23}

Traditionally, Recurrent Neural Networks (RNNs) and Long Short-Term Memory
(LSTM) networks were the standard choices for time-series forecasting due to
their ability to maintain internal state across time steps. However, since
their introduction in Natural Language Processing, Transformer models have
revolutionized time-series analysis.

Transformers use a self-attention mechanism that allows them to weigh the
importance of all past observations simultaneously. This approach inherently
captures long-range dependencies and complex multivariate interactions much
better than traditional algorithms, which often struggle to scale with
extensive historical data.

In recent years, the concept of `foundation models' has emerged for time-series
data. Models such as \textit{TimeGPT} by Nixtla~\cite{TimeGPT:01} and
\textit{TimesFM} by Google Research~\cite{TimesFM:01} have been pre-trained on
billions of data points across various domains (e.g., finance, IoT, and
weather). These models implement zero-shot forecasting, meaning they can
generate highly accurate forecasts or detect anomalies on completely unseen
data out-of-the-box, without requiring additional fine-tuning. For projects
tracking large amounts of sensor data, such as AirScout, integrating these
pre-trained transformer architectures could significantly improve the accuracy of
data validation and anomaly detection without the need to train domain-specific
models from scratch. Necessary for the training of such models is a large
amount of data, which is why we aren't using them in our project yet, but they
are definitely worth considering for future iterations of AirScout, especially
as the amount of collected data grows over time.

\begin{lstlisting}[language=Python, caption=Zero-shot Anomaly Detection using the TimeGPT Python SDK]
from nixtla import NixtlaClient
import pandas as pd

# Load time-series data from a CSV file
df = pd.read_csv('sensor_measurements.csv')

# nixtla client needs an API key, which you can get by signing up on their website
nixtla_client = NixtlaClient(api_key='YOUR_API_KEY')

# Detect anomalies out-of-the-box (Zero-shot)
anomalies_df = nixtla_client.detect_anomalies(
    df=df,
    time_col='timestamp',
    target_col='value',
    freq='15T', # 15 minute frequency
    level=99 # 99% prediction interval for the expected range
)

# Filter and display anomalies
print(anomalies_df[anomalies_df['anomaly'] == 1])
\end{lstlisting}
